\documentclass[11pt,a4paper]{article}
\usepackage[margin=2.6cm]{geometry}
\usepackage{amsmath,amssymb,booktabs}
\usepackage[colorlinks,linkcolor=blue,citecolor=blue]{hyperref}
\title{RealQM on Two Electrons:\\ one calibration, and two constructions bracketing H$_2$ and He}
\author{Claes Johnson}
\date{\today}
\begin{document}
\maketitle

\begin{abstract}
RealQM replaces the many-electron wavefunction by $N$ one-electron charge densities on
non-overlapping domains, interacting through Coulomb's law with each pair counted once. At one
electron the construction is not merely similar to the standard model---it \emph{is} the minimal
Coulomb model, with the same functional and the same minimiser, so $E(\mathrm{H}) = -1/2$ follows by
reduction. Hydrogen therefore calibrates and cannot test. The model acquires content at two
electrons, where the densities must be partitioned and must meet, and we examine the two smallest
such systems: H$_2$, whose interface lies between the kernels, and He, whose interface passes
through one. Two defensible interface treatments---strict homogeneous Neumann on a sharp partition,
and amplitude continuity enforced weakly on a diffuse one---\emph{bracket} the observed energies of
both systems from opposite sides. We argue that this interval, and not either endpoint, is what the
construction predicts, and that the width of the interval is the honest statement of its accuracy
at two electrons.
\end{abstract}

\section{The model}

Each electron $m$ carries a density $u_m^2 \ge 0$ supported on a domain $\Omega_m$, normalised to
$\int_{\Omega_m} u_m^2 = 1$, with the domains non-overlapping and filling space. The energy is
\begin{equation}
  E \;=\; \sum_m \Big[ \tfrac12 \!\int |\nabla u_m|^2 \;-\; \int K\, u_m^2 \Big]
          \;+\; \sum_{m \neq n} \tfrac12 \!\int\!\!\int \frac{u_m^2(\mathbf x)\,u_n^2(\mathbf y)}
                                                            {|\mathbf x - \mathbf y|}
          \;+\; \sum_{a<b} \frac{Z_a Z_b}{R_{ab}} ,
  \label{eq:E}
\end{equation}
with $K = \sum_a Z_a / |\mathbf x - \mathbf R_a|$ the kernel attraction, each electron pair counted
once by the $m \neq n$ rule, and no self-interaction by construction. There is no joint wavefunction
over all electrons, hence nothing to antisymmetrise and no exchange term. This is a different answer
to the same question rather than an omission---a joint antisymmetric wavefunction is not itself an
observable---and it is testable on its own terms, as below.

\paragraph{Two mechanisms, doing two different jobs.} It is worth separating them, because they are
easily run together. The Coulomb term does \emph{not} keep two electrons apart: for smooth densities
$\int\!\!\int \rho(\mathbf x)\rho(\mathbf y)/|\mathbf x - \mathbf y|$ is finite even when the two
coincide, so it permits coincidence at a finite price and cannot forbid it. What forbids it is the
\emph{partition}, which is a constraint and not a penalty: the domains are disjoint by construction,
so overlap is unavailable rather than expensive. What the Coulomb term does is resist
\emph{compression}---as two domains are pushed together $V_{ee}$ rises, and that sets where the
interface sits and how density distributes inside each domain.

The distinction matters for what follows. Between \emph{equivalent} domains across a mirror plane
the partition contributes nothing beyond what Coulomb already supplies: the sharp and continuous
constructions of \S3 agree there to $4\times 10^{-8}$. So in exactly that configuration the
construction's whole resistance to compression is the classical Coulomb repulsion of two charge
clouds, and the over-binding reported in \S5 is the measurement that this alone is not enough.
H$_2$ at its midplane and He at its split plane are both that configuration, which is why they are
the systems to examine.

\paragraph{At one electron this is the minimal Coulomb model.} For $N = 1$, \eqref{eq:E} reduces to
minimising $\tfrac12\!\int|\nabla u|^2 - \int (Z/r) u^2$ subject to $\int u^2 = 1$: the
Schr\"odinger variational problem, same functional, same minimiser, giving $-Z^2/2$ exactly. This
is worth stating in the strong form. Schr\"odinger's equation for hydrogen \emph{is} the minimal
continuum model that can be built from a Coulomb attraction---a kinetic penalty quadratic in
$\nabla u$ is the lowest-order rotation-invariant term that prevents collapse, and nothing simpler
binds at all. Hydrogen's exactness is a property of those two ingredients, not of either theory's
account of many electrons, and both theories contain exactly those ingredients at $N = 1$.
\textbf{Hydrogen therefore discriminates between no formulations whatsoever,} and nothing below
rests on it.

\paragraph{What is new, and what is missing.} Relative to one electron the construction adds exactly
two things: the densities do not overlap, and they must meet somewhere. It also \emph{subtracts}
one: it fixes what resists their approach, and the answer is set out below---the partition makes
overlap impossible while the Coulomb term alone resists compression. That is a property of the
model, stated in its own terms and checkable within it, and it is why no two-electron system should
be expected to come out exactly.

\section{One calibration, on hydrogen}

The continuum model needs no calibration. Its discrete realisation needs exactly one constant: a
grid cannot represent $1/r$, so the kernel carries a softening, $K = Z/\sqrt{r^2 + \varepsilon^2}$
with $\varepsilon = \sigma h$. We fix $\sigma$ once, by requiring the one-electron answer to be
exact. In the axisymmetric solver used below this gives
\begin{equation}
  \sigma = 0.2846 \;:\qquad E(\mathrm{H}) = -0.500227 \quad (0.045\%) ,
  \label{eq:cal}
\end{equation}
and in the three-dimensional solver $\sigma = 0.627$, read off a four-point ladder
($-0.508456$, $-0.500962$, $-0.493833$, $-0.486591$ at $\sigma = 0.4, 0.6, 0.8, 1.0$, linear in
$\sigma$ to better than a percent) and then verified at the production grid, where it returns
$E(\mathrm{H}) = -0.499901$, an error of $0.02\%$. The two values differ because the two codes
discretise differently; each is fixed by its own one-electron exactness, which is the point.

That same one-electron row carries a null test worth stating, because it costs nothing. A single
domain has no interface, so any interface measure evaluated on it must vanish; the quantity
$\int u^2 |\nabla w|\,dV$, which is the regularised form of $\oint_\Gamma u^2 dS$ used in \S5,
returns $8\times 10^{-4}$ here against $0.417$ for H$_2$. A measure that did not vanish on one
electron would be measuring something other than an interface.

Three properties of this calibration matter. It is performed on a \emph{one-electron} system, where
the model has no content of its own and the reference is exact, so no two-electron observable is
consumed. It is a single number, not a function of anything. And it is fixed before any
two-electron system is computed. Everything below follows with no further calibration.

\paragraph{What this is not.} It is not an optimisation of $\varepsilon$ against a residual. That
operation is available and must be refused: the total force residual is monotone in $\varepsilon$
and minimised by unphysical settings---hydrogen ceases to bind at $\varepsilon = 0.75h$---so a
scheme that tunes $\varepsilon$ to improve an error measure drifts away from the physics it is
meant to describe. Fixing it on a known energy is a different operation and has a unique answer.

\section{The interface, and the two constructions}

For a fixed partition, free variation of each density in \eqref{eq:E} leaves a boundary term
$\oint u_m \,\partial u_m/\partial n$, so the natural condition is homogeneous Neumann,
$\partial u/\partial n = 0$, on each domain's own boundary---nothing is imposed. Solved that way the
amplitudes at the interface generically \emph{disagree}, and the interface is then placed where they
agree. Continuity is therefore not a second assumption about $\Psi$; it is the equation that
determines $\Gamma$. Three conditions, one extra unknown.

Two implementations each supply one half of that scheme and take the other half from elsewhere.

\paragraph{(A) Sharp partition, strict Neumann.} Labelled cells, with the kinetic flux through any
face joining cells of different ownership \emph{omitted at assembly}. Omitting a face \emph{is} the
homogeneous-Neumann condition---zero flux implemented by not adding a flux, with no extra equation.
This imposes the condition exactly but requires $\Gamma$ to be known, which for H$_2$ and He is
supplied by symmetry. One trap is worth recording: taking a \emph{submatrix} of the full stiffness
instead retains the diagonal contributions of faces leading into the excluded region and so imposes
Dirichlet, pushing density off the boundary and under-binding. Measured on a one-electron problem
outside a core, $3.977$ against $6.177\,\mathrm{eV}$; the latter matches a one-dimensional radial
solve to $2\%$.

\paragraph{(B) Diffuse partition, continuity enforced weakly.} An occupancy $w_m$ evolves on the
local amplitude mismatch $c = u_m - \max_{n\neq m} u_n$ and converges where that mismatch vanishes,
so this construction genuinely \emph{finds} $\Gamma$---the only one of the two that does. Its
Neumann condition is only natural, not imposed: varying $\tfrac12\int w|\nabla u|^2$ makes
$w\,\partial u/\partial n = 0$ stationary, but on a finite grid a face straddling the interface
carries half flux rather than none. That half-flux is an extra kinetic cost sitting exactly at the
interface, and it is a first-order discretisation error, not a structural one: H$_2$'s bond error
falls $34.9 \to 28.3 \to 21.4\%$ at $h = 0.1395, 0.10, 0.08$ and is still falling.

\paragraph{The two are not a choice between theories.} They are two incomplete implementations of
one: (A) imposes the condition and is told where $\Gamma$ is; (B) finds $\Gamma$ and imposes the
condition weakly. Combining them---binary faces on a \emph{diffuse} occupancy---does not work and
should not be attempted: it gives H$_2$ a binding energy of $+110\%$ and a dipole of
$0.108$--$0.187\,\mathrm{D}$ where reflection symmetry requires identically zero, because closed
faces make an amplitude jump at the interface free. The diagnostic is direct: the kinetic energy of
the assembled field $\Psi = \sum_m w_m u_m$, which charges that jump, exceeds the per-domain sum by
a factor of $8.07$, so the discrete solution is not in $H^1$.

\section{H$_2$ and He are the two minimal two-electron systems}

They are not a repetition. In H$_2$ the interface lies \emph{between} the two kernels, in the tail
of both densities; in He it passes \emph{through} the single kernel, cutting the density maximum.
Measured on the same construction, $\int_\Gamma u^2 \,dS = 0.496$ for H$_2$ against $1.681$ for
He---a factor of $3.4$. Any interface treatment is therefore tested far more severely by He, and an
error in it cannot be hidden there.

One asymmetry must be declared. H$_2$'s interface position is fixed by reflection symmetry about the
midplane, and is the same in both constructions. He's is not: nothing in a spherical two-electron
atom selects a plane, so construction (A) imposes one by hand and the result carries that
arbitrary orientation. H$_2$ is in this respect the cleaner system and He the more demanding one.

\paragraph{And H$_2$ cannot settle the interface question.} At a symmetry plane a mirror-symmetric
solution has $u_1 = u_2$ \emph{and}, by reflection, $\partial u/\partial n = 0$. The two conditions
coincide, and the two constructions agree there to $4\times 10^{-8}$. Settling which is right
requires an asymmetric interface, which neither of these systems provides.

\section{The bracket}

\begin{table}[t]
\centering
\begin{tabular}{lrrr}
\toprule
 & (A) sharp, strict Neumann & observed & (B) diffuse, weak \\
\midrule
$E(\mathrm{H})$   & $-0.500227$ & $-0.5$      & calibrated \\
$E(\mathrm{H_2})$ & $-1.278929$ & $-1.174470$ & $-1.132147$ \\
$E(\mathrm{He})$  & $-3.026749$ & $-2.903724$ & $\approx -2.60$ \\
\bottomrule
\end{tabular}
\caption{The two constructions straddle both observed energies. (A) at its own calibration
\eqref{eq:cal}, H$_2$ at its own energy minimum. (B) at $\sigma = 1$, which is \emph{not} its
calibrated value---rows at $\sigma = 0.627$ are being measured and will move column (B) toward the
centre, narrowing the interval. The observed H$_2$ and He figures are themselves computations in
the standard theory, not measurements; see the caveat below.}
\label{tab:bracket}
\end{table}

Table~\ref{tab:bracket} is the paper's result. Construction (A) over-binds both systems; (B)
under-binds both; the observed values lie between. In relative terms (A) gives $E_{\text{bind}}$ for
H$_2$ as $+59.6\%$ at its own minimum while placing the bond within $-3.7\%$, and He as $4.2\%$ too
deep; (B) under-binds H$_2$ and gives He some $10\%$ too shallow.

\paragraph{The interval is the prediction.} Neither endpoint should be quoted as the construction's
answer. Both are defensible readings of the same model differing only in how the interface is
treated, and selecting the one nearer experiment would be fitting a discrete choice---the more so
because an intermediate treatment can be constructed at will. What the model predicts at two
electrons is an interval, and its width is the honest measure of the accuracy the interface
treatment currently permits. The useful question is not which endpoint is right but whether the
interval narrows under refinement, and for (B) it demonstrably does: the half-flux error is
first-order in $h$.

\paragraph{A one-parameter repair exists and we do not adopt it.} The surface energy
$(\beta/2)\oint_\Gamma \Psi^2 dS$, whose natural condition is the Robin condition
$\partial u/\partial n + \beta u = 0$, interpolates between (A) at $\beta = 0$ and a Dirichlet
limit. It is the repulsion between \emph{equivalent} domains across a mirror plane that territorial
exclusion cannot supply, so it acts exactly where the over-binding of (A) lives. Calibrated on
H$_2$'s binding energy it gives $\beta = 0.209$, i.e.\ $\beta^{-1} = 4.78\,a_0$, a plausible length;
and carried unchanged to He it improves that system from $-4.24\%$ to $+1.62\%$, with He's own
preferred value $0.151$---agreement to a factor $1.4$ across two systems with very different
interface geometry. We nonetheless report $\beta = 0$, for one reason: $\beta$ cannot deliver energy
and geometry together. At H$_2$'s own minimum it trades $+59.6\% \to +18\%$ on binding for
$-3.7\% \to +7.6\%$ on the bond, with a balanced point near $\beta = 0.25$ at roughly $\pm 10\%$ on
both. One parameter cannot satisfy both conditions, and adopting it would spend the construction's
only real asset---that nothing is adjusted---for an improvement in one observable at the expense of
the other.

\section{What this establishes, and what it does not}

\paragraph{Not that the model is accurate at two electrons.} It is not: the interval of
Table~\ref{tab:bracket} is tens of percent wide in the energy. That is the expected size of a
declared property of the construction---territorial exclusion is all there is, and it vanishes
between equivalent domains across a mirror plane---rather than a numerical failing to be driven
down, and the
testable claim is that the error is \emph{systematic}: one sign, one scale, and cancelling in
differences taken at fixed composition.

\paragraph{A caveat on the references.} The figure $-1.174470$ for H$_2$ is a Kolos--Wolniewicz
\emph{computation}, in the very theory whose many-electron content this construction declines, and
the total electronic energy is in any case not an observable. The observables here are the
dissociation energy, the bond length, and a dipole that must vanish by symmetry. We report total
energies because they are what the two constructions differ in, not because they are measurable.

\paragraph{What the dipole buys.} For a symmetric system the computed dipole must be identically
zero, so any nonzero value is a pure error measurement needing no reference of any kind. It is the
one check in this paper that cannot be argued with, and it is what disqualified the combined
construction of \S3 and would disqualify any future variant that improved the energy while losing
the symmetry.

\section{HeH$^+$: the asymmetric case}

Both systems above have interfaces fixed by symmetry, where the two constructions provably coincide,
so neither can settle which formulation is right. The smallest two-electron system that can is
HeH$^+$: two electrons on kernels of charge $2$ and $1$, net charge $+1$, equilibrium separation
$R_e = 1.4632\,a_0$. Its interface lies at no symmetry plane, so amplitude matching and zero flux
are genuinely different conditions there and the $\sim$\,5-fold discrepancy they show at a displaced
interface should appear in an observable.

The construction is the natural one and needs nothing new: two domains, one electron each, one
carrying a $Z = 2$ kernel and the other $Z = 1$. There is no pseudopotential, no core, no split
ambiguity, and $r_c = 0$ throughout.

\paragraph{One confound must be declared, and it is not small.} The two kernels have
\emph{different} charges, so a single softening $\sigma$ cannot be calibrated for both: the value
fixed on hydrogen is not the value a $Z = 2$ cusp requires, and the sharper cusp is the less well
resolved. An error in HeH$^+$ is therefore attributable either to the interface asymmetry---the
thing we want to measure---or to the $Z$ mismatch, and the two are not separated by this system
alone. The honest reading is that HeH$^+$ tests the formulation only to the extent that the $\sigma$
transfer error is smaller than the effect sought, which has not been established.

\paragraph{The clean alternative, with equal kernels.} The confound disappears if the asymmetry is
put in the interface rather than in the nuclei: H$_2$ at fixed $D$ with the dividing plane
\emph{displaced} off-centre, each domain still carrying unit charge. Both kernels are then $Z = 1$,
one $\sigma$ serves both, and the only asymmetry is the one under test. That measurement exists.
Writing $\delta E = \tfrac12 k z_i^2$ for a displacement $z_i$,
\begin{center}
\begin{tabular}{lccc}
\toprule
$h$ & $k$, continuity enforced & $k$, strict per-domain & ratio \\
\midrule
$0.250$ & $0.7359$ & $0.1296$ & $5.68$ \\
$0.125$ & $0.7394$ & $0.1472$ & $5.02$ \\
\bottomrule
\end{tabular}
\end{center}
The continuity stiffness is converged---$0.5\%$ across a halving---while the strict one is not,
climbing $13.6\%$ and still moving. That asymmetry is the sharper statement: the strict construction
charges for a \emph{discontinuity}, a mesh-scale object, so it has no well-defined interface
stiffness in the continuum limit, whereas the continuity construction does. For an interface that
must \emph{move}---which is every molecule without a symmetry plane---that is a difference in kind
rather than in accuracy.

\paragraph{So the programme is: HeH$^+$ for the observable, displaced H$_2$ for the mechanism.}
HeH$^+$ gives a real system with a real reference and a genuinely asymmetric interface, at the cost
of the $\sigma$ confound; displaced H$_2$ isolates the mechanism exactly but is not an observable
configuration. Neither alone decides it; together they do, and the prediction to test is that the
construction which converges in interface stiffness is the one whose geometry converges in
HeH$^+$.

\end{document}
